\begin{name}
	{\tenchude}
	{TOÁN 10}
	{LỚP TOÁN THẦY PHÁT}
	{Cuộc sống như chiếc gương, bạn chỉ nhận được kết quả tốt đẹp nhất khi mỉm cười với nó.}
\end{name}
\Opensolutionfile{ansbook}[ans/ansbookDe1]
\TN
\Opensolutionfile{ans}[ans/ansDe1-TN1]
\begin{ex}%[0D8N1-1]
Một lớp học có $25$ học sinh nam và $20$ học sinh nữ. Hỏi có bao nhiêu cách chọn ra một học sinh trong lớp học này đi dự trại hè của trường?
\choice
{$25$}
{$20$}
{\True$45$}
{$500$}
\loigiai
{
Áp dụng quy tắc cộng:
Số cách chọn ra một học sinh trong lớp học này đi dự trại hè của trường là 25 + 20 = 45.
}
\end{ex}

\begin{ex}%[0D8N2-4]
Trong kì thi TN THPT Quốc gia năm $2023$ tại một điểm thi có $5$ sinh viên tình nguyện được phân công trực hướng dẫn thí sinh ở $5$ vị trí khác nhau. Yêu cầu mỗi vị trí có đúng $1$ sinh viên. Hỏi có bao nhiêu cách phân công vị trí trực cho $5$ sinh viên đó?
\choice
{\True $120$}
{$25$}
{$10$}
{$24$}
\loigiai{
Số cách phân công $5!=120$ cách.
}
\end{ex}

\begin{ex}%[BG - 10 New - 4in1,Phú Thạch]%[0D8N3-1]
Có bao nhiêu số hạng trong khai triển nhị thức $(2023x+2024)^4$?
\choice
{$4$}
{$6$}
{\True $5$}
{$3$}
\loigiai
{
Số số hạng trong khai triển nhị thức $(2023x+2024)^4$ là $4+1=5$.
}
\end{ex}

\begin{ex}%[0D6H1-1]
	Kết quả đo chiều dài của một cây cầu được ghi là $152m\pm 0,2m$. Tìm sai số tương đối của phép đo chiều dài cây cầu đó
	\choice
	{\True $\delta_a<0,1316\%$}
	{$\delta_a<1,316\%$}
	{$\delta_a=0,1316\%$}
	{$\delta_a>0,1316\%$}
	\loigiai{
		Sai số tương đối $\delta_a\le \dfrac{0,2}{152}=0,001315789\approx 0,1316\%$.
	}
\end{ex}

\begin{ex}%[VN-MT-9]%[ 0-TK-HK1-KN-2-2425 , Trần Đức Thắng]%[0D6H4-2]
	Mẫu số liệu hình bên cho biết cân nặng của $10$ trẻ sơ sinh (đơn vị: kg).
	\begin{center}
		\begin{tabular}{|c|c|c|c|c|c|c|c|c|c|}
			\hline
			$2{,}977$ & $3{,}155$ & $3{,}920$ & $4{,}236$ & $3{,}387$ & $2{,}593$ & $3{,}270$ & $3{,}412$ & $3{,}813$ & $4{,}042$ \\
			\hline
		\end{tabular}
	\end{center}
	Khoảng biến thiên (đơn vị: kg) của mẫu số liệu trên bằng
	\choice
	{$0{,}384$}
	{$0{,}194$}
	{\True $1{,}643$}
	{$3$}
	\loigiai{
	Trước hết, ta sẽ sắp xếp mẫu số liệu theo thứ tự không giảm
	\begin{center}
		\begin{tabular}{|c|c|c|c|c|c|c|c|c|c|}
			\hline
			$ 2{,}593$ & $ 2{,}977$ & $ 3{,}155 $ & $ 3{,}270 $ & $3{,}387$ & $ 3{,}412 $ & $3{,}813 $ & $ 3{,}920 $ & $4{,}042$ & $ 4{,}236$ \\
			\hline
		\end{tabular}
	\end{center}
	Khoảng biến thiên là $R=4{,}236-2{,}593=1{,}643$ (kg).
	}
\end{ex}

\begin{ex}%[0D0N1-2]
	Gieo hai đồng tiền một lần. Kí hiệu $S$, $N$ lần lượt để chỉ đồng tiền lật sấp, lật ngửa. Xác định biến cố $M$: \lq\lq Hai đồng tiền xuất hiện hai mặt không giống nhau\rq\rq.
	\choice
	{$M = \left\{NN; SS\right\}$}
	{\True $M = \left\{NS; SN\right\}$}
	{$M = \left\{NS; NN\right\}$}
	{$M = \left\{SS; SN\right\}$}
	\loigiai{
		Biến cố $M$: \lq\lq Hai đồng tiền xuất hiện hai mặt không giống nhau\rq\rq\,là $M = \left\{NS; SN\right\}$.
	}
\end{ex}

\begin{ex}%[0D0N1-3]
	Chọn ngẫu nhiên một cuốn truyện từ $27$ cuốn truyện cổ tích khác nhau và $19$ cuốn tiểu thuyết trinh thám khác nhau. Số phần tử của không gian mẫu là
	\choice
	{$50$}
	{$513$}
	{$8$}
	{\True$46$}
	\loigiai{
		Số phần tử của không gian mẫu là $n(\Omega)=27+19=46$.

	}
\end{ex}

\begin{ex}%[0H9N3-1]%[tex hóa đề CK2 - 2025 - Huyên Nguyễn]
	Trong mặt phẳng tọa độ $Oxy$, đường thẳng $\Delta$: $2x-y+2025=0$ có một vectơ pháp tuyến là
	\choice
	{$\overrightarrow{n}=(1;2)$}
	{\True$\overrightarrow{n}=(4;-2)$}
	{$\overrightarrow{n}=(2;1)$}
	{$\overrightarrow{n}=(-2;-1)$}
	\loigiai{Đường thẳng  $\Delta$: $2x-y+2025=0$ có một vectơ pháp tuyến là $\overrightarrow{n}=(4;-2)$.
	}
\end{ex}

\begin{ex}
	Trong mặt phẳng tọa độ $Oxy$, cho đường thẳng $d_1$ có vectơ chỉ phương $\vec{u_1}$ và đường thẳng $d_2$ có vectơ chỉ phương $\vec{u_2}$. Góc giữa hai đường thẳng $d_1$ và $d_2$ được tính bằng công thức
	\choice
	{$\cot (d_1,d_2)=\dfrac{|\vec{u_1}\cdot \vec{u_2}|}{|\vec{u_1}||\vec{u_2}|}$}
	{\True $\cos (d_1,d_2)=\dfrac{|\vec{u_1}\cdot \vec{u_2}|}{|\vec{u_1}||\vec{u_2}|}$}
	{$\sin (d_1,d_2)=\dfrac{|\vec{u_1}\cdot \vec{u_2}|}{|\vec{u_1}||\vec{u_2}|}$}
	{$\tan (d_1,d_2)=\dfrac{|\vec{u_1}\cdot \vec{u_2}|}{|\vec{u_1}||\vec{u_2}|}$}
	\loigiai{
		Gọi $\alpha$ là góc giữa hai đường thẳng $d_1$ và $d_2$.\\
		Ta có $\cos \alpha = \dfrac{|\vec{u_1}\cdot \vec{u_2}|}{|\vec{u_1}||\vec{u_2}|}$.
	}
\end{ex}

\begin{ex}%[10-11EX-GK2 24-25]%[Lương Như Quỳnh]%[0H9N3-3]
	Trong mặt phẳng $Oxy$, đường thẳng $d\colon x-2y-1=0$ song song với đường thẳng có phương trình nào sau đây?
	\choice
	{$-x+2y+1=0$}
	{\True $-2x+4y-1=0$}
	{$x+2y+1=0$}
	{$2x-y=0$}
	\loigiai{
		Xét đường thẳng $d\colon x-2y-1=0$ và đường thẳng $\Delta \colon -2x+4y-1=0$.\\
		Ta có $\dfrac{1}{-2}=\dfrac{-2}{4}\neq \dfrac{-1}{-1} $.\\
		Suy ra $d\parallel \Delta$.
	}
\end{ex}

\begin{ex}%[0H9N5-1]
	Một tiêu điểm của elip $(E)\colon \dfrac{x^2}{5}+\dfrac{y^2}{4}=1$ là
	\choice
	{$F_2=(0;1)$}
	{\True $F_2=(1;0)$}
	{$F_1=(-3;0)$}
	{$F_1=(-1;-2)$}
	\loigiai{
		Ta có $a^2=5$, $b^2=4 \Rightarrow c^2=a^2-b^2=1 \Rightarrow c=1 \Rightarrow F_1=(-1;0), F_2=(1;0)$.
	}
\end{ex}

\begin{ex}%[0H9N5-5]
	Phương trình nào sau đây là phương trình chính tắc của hypebol?
	\choice
	{$x^2+\dfrac{y^2}{3^2}=1$}
	{$\dfrac{x^2}{16}-y^2=-1$}
	{$\dfrac{x^2}{25}-\dfrac{y^2}{9}=-1$}
	{\True $x^2-\dfrac{y^2}{2}=1$}
	\loigiai{
		Phương trình
		$x^2-\dfrac{y^2}{2}=1 \Leftrightarrow \dfrac{x^2}{1}-\dfrac{y^2}{2}=1$ là phương trình chính tắc của hypebol.
	}
\end{ex}
\Closesolutionfile{ans}

\TNTF
\Opensolutionfile{ans}[ans/ansDe1-TN2]
\begin{ex}%[Dự án EX-10-11-Chuẩn hóa]%[Hoàng Thanh Phương]%[0D6H4-3]
	Mẫu số liệu sau cho biết số ghế trống của một rạp chiếu phim trong 11 ngày
	\begin{center}
	\begin{tabular}{|l|l|l|l|l|l|l|l|l|l|l|}
	\hline 0 & 7 & 8 & 22 & 20 & 15 & 18 & 19 & 13 & 11 & 39 \\
	\hline
	\end{tabular}
	\end{center}
	Các khẳng định sau là đúng hay sai?
	\choiceTF
	{\True Khoảng tứ phân vị của mẫu số liệu là $12$}
	{Không có giá trị bất thường trong mẫu số liệu nêu trên}
	{Quy tròn số trung bình của mẫu số liệu với độ chính xác $d=0.5$ là $a=15$}
	{\True Sai số tương đối của số quy tròn $a$ xấp xỉ $2,27\%$}
	\loigiai{
	Sắp xếp mẫu số liệu theo thứ tự không giảm ta được
	\[0\ 7\ 8\ 11\ 13\ 15\ 18\ 19\ 20\ 22\ 39\]
	Ta thấy $Q_1=8, Q_2=15, Q_3=20$.\\
	\begin{itemchoice}
	\itemch Khoảng tứ phân vị $\Delta_Q=20-8=12$.
	\itemch Ta có $Q_1-\dfrac{3}{2}\Delta_Q=-10$, $Q_3+\dfrac{3}{2}\Delta_Q=38$.\\
	Vậy $39$ là giá trị bất thường của mẫu số liệu trên.
	\itemch Số trung bình cộng của mẫu số liệu là
	\[\overline{x}=\dfrac{0+7+8+11+13+15+18+19+20+22+39}{11}=\dfrac{172}{11}\approx 15,6363.\]
	Vậy quy tròn số trung bình cộng với độ chính xác $d=0,5$ là $16$.
	\itemch Sai số tương đối của số quy tròn $a$ là
	\[\delta_a=\dfrac{|a-\overline{x}|}{a}=\dfrac{|16-\dfrac{172}{11}|}{16} \approx 2,27\%.\]
	\end{itemchoice}
	}
	\end{ex}

\begin{ex}%[0H9H3-2]
	Cho hai điểm $A(3;-3)$, $B(-1;-5)$ và đường thẳng $(d)\colon 4x-3y-2=0$.
	\choiceTF
	{\True Một vectơ chỉ phương của đường thẳng $d$ là $\overrightarrow{u_d}=(3;4)$}
	{Đường thẳng đi qua điểm $A$ và vuông góc với $(d)$ có phương trình $4x+3y=3$}
	{Đường tròn $(C)$ đường kính $AB$ có phương trình là ${(x-1)^2+(y+4)^2=10}$}
	{\True Đường tròn $(C)$ cắt đường thẳng $(d)$ tại hai điểm phân biệt}
	\loigiai{
		\begin{itemchoice}
			\itemch  \textbf{Đúng}.\\  Vectơ pháp tuyến của  đường thẳng $d$ là $\overrightarrow{n_d}=(4;-3)$ suy ra vectơ chỉ phương  $\overrightarrow{u_d}=(3;4)$.
			\itemch \textbf{Sai}.\\  Đường thẳng vuông góc với đường thẳng $(d)$ nên có dạng: $\Delta\colon 3x+4y+m=0$.\\
			$\Delta$ đi qua $A$ nên thay toạ độ $A$ vào phương trình đường thẳng $\Delta$ ta có \[3\cdot3+4\cdot(-3)+m=0\Leftrightarrow m=3.\]
			Vậy phương trình đường thẳng cần tìm là $\Delta\colon 3x+4y+3=0$.
			\itemch \textbf{Sai}.\\
			Phương trình đường tròn $(C)$ có tâm $I(1;-4)$ và bán kính $R=\dfrac{AB}{2}=\sqrt{5}$.\\
			Vậy phương trình đường tròn $(C)$ là $(x-1)^2+(y+4)^2=5$
			\itemch \textbf{Đúng}.\\
			Ta có $\mathrm{d}(I,(d))=\dfrac{|4\cdot1-3\cdot(-4)-2|}{\sqrt{4^2+(-3)^2}}=\dfrac{14}{5}>\sqrt{5}$.\\
			Vậy đường tròn $(C)$ cắt đường thẳng $(d)$ tại hai điểm phân biệt.
		\end{itemchoice}
	}

\end{ex}
\Closesolutionfile{ans}

\TNSA
\Opensolutionfile{ans}[ans/ansDe1-TN3]
\begin{ex}%[H]giảng 10 New - 4in1, Nguyễn Xuân Bảo]%[0D8H1-2]
Số $2531250000$ có bao nhiêu ước số tự nhiên chia hết cho $5$?
\shortans{$160$}
\loigiai{
Ta có $253125000=2^3\cdot 3^4\cdot 5^8$ nên mỗi ước số tự nhiên của số đã cho đều có dạng $2^m\cdot 3^n\cdot 5^p$ trong đó $m$, $n$, $p\in\mathbb{N}$ sao cho $0\le m\le 3$, $0\le n\le 4$, $0\le p\le 8$.
\begin{itemize}
\item Gọi $A$ là tập hợp các ước số tự nhiên của $253125000$, $n(A)=4\cdot 5 \cdot 9=180$.
\item  Gọi $B$ là tập hợp các ước số tự nhiên không chia hết cho $5$ của $253125000$, $n(B)=4\cdot 5 =20$.
\end{itemize}
Vậy số ước số tự nhiên chia hết cho $5$ của $2531250000$ là $180-20=160$.
}
\end{ex}

\begin{ex}%[Chuẩn hóa M25, Nguyễn Văn Nay]%[0D8V2-3]
Có bao nhiêu số tự nhiên có $5$ chữ số đôi một khác nhau, mà trong mỗi số đều có ba chữ số $0$, $1$, $2$?
\shortans{$2016$}
\loigiai{
Gọi $\overline{abcde}$ là số cần lập.
\begin{itemize}
\item Sắp xếp $3$ vị trí trong $5$ vị trí cho ba chữ số $0$, $1$, $2$ có $\mathrm{A}_5^3=60$ cách.
\item Sắp xếp $2$ vị trí còn lại có $\mathrm{A}_7^2=42$ cách.
\item Do đó có $60\cdot 42=2520$ số.
\end{itemize}
Gọi $\overline{0bcde}$ là số mà chữ số $0$ đứng đầu.
\begin{itemize}
\item Sắp xếp $2$ vị trí trong $4$ vị trí cho ba chữ số $1$, $2$ có $\mathrm{A}_4^2=12$ cách.
\item Sắp xếp $2$ vị trí còn lại có $\mathrm{A}_7^2=42$ cách.
\item Do đó có $12\cdot 42=504$ số.
\end{itemize}
Vậy có $2520-504=2016$ số.
}
\end{ex}

\begin{ex}%[0H9V5-0]%[tex hóa đề ck2 - form 2025 - đợt 2- Duong Xuan Loi]
\immini{
Một nhà vòm chứa máy bay có mặt cắt hình nửa elip cao $5$ m, rộng $20$ m. Khoảng cách theo phương thẳng đứng từ một điểm cách chân tường $5$ m lên đến nóc nhà vòm bằng $\dfrac{a\sqrt{b}}{c}$ với $a$, $b$, $c$ là các số nguyên dương. Tính giá trị biểu thức $T=a+2b-c$.
}{
\begin{tikzpicture}[>=stealth,line join=round,line cap=round,font=\footnotesize,scale=.45]
%mái vòm
\draw[green!50!black] (8,0) arc(0:180:8 and 7.8);
\foreach \i in {1,2,3,4,5,6,7,8,9,10,11}
\draw[green!50!black] (0,0)--({15*(\i)}:8 and 7.8);
\fill[draw=green!50!black,white] (7.6,0)arc(0:180:7.6 and 7.5)--cycle;
\draw[green!50!black] (7.6,0) arc(0:180:7.6 and 7.5);

\draw[green!50!black] (6.7,0.8) arc(0:180:6.7 and 6.3);
\foreach \i in {1,2,3,4,5,6,7,8,9,10,11}
\draw[green!50!black] (0,0.8)--($(0,0.8)+({15*(\i)}:6.7 and 6.3)$);
\fill[draw=green!50!black,white] (6.4,0.8)arc(0:180:6.4 and 6.1)--cycle;
\draw[green!50!black] (6.4,0.8) arc(0:180:6.4 and 6.1);

\draw[green!50!black] (5.4,1.6) arc(0:180:5.4 and 4.8);
\foreach \i in {1,2,3,4,5,6,7,8,9,10,11}
\draw[green!50!black] (0,1.6)--($(0,1.6)+({15*(\i)}:5.4 and 4.8)$);
\fill[draw=green!50!black,white] (5.15,1.6)arc(0:180:5.15 and 4.6)--cycle;
\draw[green!50!black] (5.15,1.6) arc(0:180:5.15 and 4.6);

\foreach \i in {1,2,3,4,5,6,7,8}
\draw[green!50!black] (0,3.5)--({20*(\i)}:8 and 7.8);
\fill[draw=green!50!black,white] (4.1,2.4)arc(0:180:4.1 and 3.3)--cycle;
\draw[green!50!black] (4.1,2.4) arc(0:180:4.1 and 3.3);
\foreach \i in {1,2,3,4,5,6,7,8,9,10,11}
\draw[green!50!black] (0,2.4)--($(0,2.4)+({15*(\i)}:4.1 and 3.3)$);
\fill[draw=green!50!black,white] (3.85,2.4)arc(0:180:3.85 and 3.1)--cycle;
\draw[green!50!black] (3.85,2.4) arc(0:180:3.85 and 3.1);
\fill[middle color=cyan,opacity=0.05] (8,0) arc(0:180:8 and 7.8);
%cửa sau
\fill[draw=green!50!black,bottom color=cyan!20!white,top color=white] (3.85,2.4)arc(0:180:3.85 and 3.1)--cycle;
\draw[color=green!50!black] ($(0,2.4)+(60:3.85 and 3.1)$)|-(0,2.4) coordinate[pos=0.1](A);
\draw[color=green!50!black] ($(0,2.4)+(120:3.85 and 3.1)$)|-(0,2.4) coordinate[pos=0.1](B);
\draw[color=green!50!black] ($(0,2.4)+(60:3.85 and 3.1)$)|-($(0,2.4)+(30:3.85 and 3.1)$) ($(0,2.4)+(120:3.85 and 3.1)$)|-($(0,2.4)+(150:3.85 and 3.1)$) (A)--(B);
%sàn nhà
\fill[bottom color=gray!40!white,top color=white] (8,0)--(4.1,2.4)--(-4.1,2.4)--(-8,0)--cycle;
%bánh xe bên phải
\fill[color=black] (1,1.4)--(1.5,0.9)--++(0.1,0)--(1.1,1.4)--cycle;
\fill[color=black] (1.5,0.8)--++(0,0.7)--++(0.1,0)--++(0,-0.7)--++(0.3,0)--++(0,-0.5)--++(-0.2,0)--++(0,0.4)--++(-0.3,0)--++(0,-0.4)--++(-0.2,0)--++(0,0.5)--++(0.3,0)--cycle;
\fill[left color=gray] (2.6,1) circle(0.5 cm);
\fill[right color=black] (2.6,1) circle(0.4 cm);
\fill[left color=gray] (3.7,1.23) circle(0.4 cm);
\fill[right color=black] (3.7,1.23) circle(0.3 cm);
%bánh xe bên trái
\fill[xscale=-1,color=black] (1,1.4)--(1.5,0.9)--++(0.1,0)--(1.1,1.4)--cycle;
\fill[xscale=-1,color=black] (1.5,0.8)--++(0,0.7)--++(0.1,0)--++(0,-0.7)--++(0.3,0)--++(0,-0.5)--++(-0.2,0)--++(0,0.4)--++(-0.3,0)--++(0,-0.4)--++(-0.2,0)--++(0,0.5)--++(0.3,0)--cycle;
\fill[xscale=-1,left color=gray] (2.6,1) circle(0.5 cm);
\fill[xscale=-1,right color=black] (2.6,1) circle(0.4 cm);
\fill[xscale=-1,left color=gray] (3.7,1.23) circle(0.4 cm);
\fill[xscale=-1,right color=black] (3.7,1.23) circle(0.3 cm);
%cánh phải dưới
\fill[left color=gray,right color=cyan!50!black] (0.75,1.5)--(6,1.9)--++(0.5,0.35)--++(-0.2,0)--(6,2)--(0.7,1.9)--cycle;
\fill[left color=yellow!40!white, right color=cyan!50!black] (0.75,1.5)--(6,1.9)--($(0,1.3)+(-5:0.8 and 0.5)$)arc(0:-94:0.75 and 0.4)--($(0,1.5)+(-90:0.75 and 0.5)$) arc(-90:0:0.75 and 0.5)--cycle;
%cánh trái dưới
\fill[xscale=-1,right color=gray,left color=cyan!50!black] (0.75,1.5)--(6,1.9)--++(0.5,0.35)--++(-0.2,0)--(6,2)--(0.7,1.9)--cycle;
\fill[xscale=-1,right color=yellow!40!white, left color=cyan!50!black] (0.75,1.5)--(6,1.9)--($(0,1.3)+(-5:0.8 and 0.5)$)arc(0:-94:0.75 and 0.4)--($(0,1.5)+(-90:0.75 and 0.5)$) arc(-90:0:0.75 and 0.5)--cycle;
%cánh phải trên
\fill[left color=gray,right color=cyan!50!black] (0.7,2)--(3.1,2.15)--++(0.05,0.05)--($(0,2)+(15:0.7 and 0.5)$)--cycle;
%cánh trái trên
\fill[xscale=-1,right color=gray,left color=cyan!50!black] (0.7,2)--(3.1,2.15)--++(0.05,0.05)--($(0,2)+(15:0.7 and 0.5)$)--cycle;
%đuôi máy bay
\fill[bottom color=yellow!40!white,top color=cyan!40!black] ($(0,2)+(80:0.7 and 0.5)$)--(0,5)--($(0,2)+(100:0.7 and 0.5)$)--cycle;
%đầu máy bay
\fill[bottom color=black!70, top color=cyan] (0.7,2) arc(0:180:0.7 and 0.5)--(-0.75,1.5) arc(180:0:0.75 and 0.2)--cycle;
\fill[top color=black!70, bottom color=cyan!50!white](-0.75,1.5) arc(-180:0:0.75 and 0.5) arc(0:180:0.75 and 0.2)--cycle;
\fill[color=black] (0.75,1.5)--(0.7,2)--(0,1.9)--(-0.7,2)--(-0.75,1.5) arc(180:0:0.75 and 0.2)--cycle;
\fill[pattern color=gray, pattern=vertical lines] (0.75,1.5)--(0.7,2)--(0,1.9)--(-0.7,2)--(-0.75,1.5) arc(180:0:0.75 and 0.2)--cycle;
%bánh xe giữa
\fill[color=black] ($(0,1.3)+(-90:0.75 and 0.4)$)--(0.1,0.6)--++(0.1,0)--++(0,-0.3)--++(-0.1,0)--++(0,0.25)--($(0,1.25)+(-90:0.75 and 0.4)$)--cycle;
\fill[xscale=-1,color=black] ($(0,1.3)+(-90:0.75 and 0.4)$)--(0.1,0.6)--++(0.1,0)--++(0,-0.3)--++(-0.1,0)--++(0,0.25)--($(0,1.25)+(-90:0.75 and 0.4)$)--cycle;
\end{tikzpicture}
}
\shortans[0]{$9$}
\loigiai{
Chọn hệ trục tọa độ $Oxy$ với gốc tọa độ tại tâm đáy nhà vòm, trục tung thẳng đứng.
\begin{center}
\begin{tikzpicture}[scale=0.5, font=\footnotesize, line join=round, line cap=round,>=stealth]
\def\xmin{-14} \def\xmax{14}
\def\ymin{-7} \def\ymax{7}
\draw[->] (\xmin,0)--(\xmax,0) node [below]{$x$};
\draw[->] (0,\ymin)--(0,\ymax) node [left]{$y$};
\node at (0,0) [below left]{$O$};
\clip (\xmin+0.1,\ymin+0.1) rectangle (\xmax-0.1,\ymax-0.1);
\draw (0,0) ellipse (10cm and 5cm);
\fill (10,0) circle (2.0pt) node[above right]{$C(10;0)$}
(-10,0) circle (2.0pt) node[above left]{$B(-10;0)$}
(0,5) circle (2.0pt) node[above right]{$A(0;5)$};
\end{tikzpicture}
\end{center}
Nhà vòm có dạng nửa elip nên có phương trình chính tắc của elip là $\dfrac{x^2}{a^2}+\dfrac{y^2}{b^2}=1$ ($a, b>0$).\\
Ta có chiều cao của nhà vòm là $5$ m nên $OA=h=5$, chiều rộng của nhà vòm là $20$ m nên $BC=2OB=20$. Suy ra $OB=10$.\\
Ta có tọa độ các điểm: $C(10; 0)$ và $A(0; 5)$. Thay hai điểm này vào phương trình chính tắc, ta có\\
$\heva{&\dfrac{10^2}{a^2}+\dfrac{0^2}{b^2}=1\\&\dfrac{0^2}{a^2}+\dfrac{5^2}{b^2}=1}\Leftrightarrow\heva{&a=10\\&b=5\cdot.}$ \\
Suy ra phương trình miêu tả hình dáng nhà vòm là $\dfrac{x^2}{100}+\dfrac{y^2}{25}=1$.\\
Điểm cách chân tường $5$ m tương ứng cách tâm $5$ m (vì từ tâm vòm đến tường là $10$ m).\\
Thay $x=5$ vào phương trình $\dfrac{x^2}{100}+\dfrac{y^2}{25}=1$, ta tìm được $y=\dfrac{5\sqrt{3}}{2}$.\\
Khi đó: $\heva{&a=5\\&b=3\\&c=2}\Rightarrow T=a+2b-c=5+2\cdot 3-2=9$.}
\end{ex}

\begin{ex}%[0H9V4-2]
	Trong mặt phẳng $Oxy$, gọi $(C)$ là đường tròn đi qua hai điểm $A( 1;1)$, $B(5;3)$ và có tâm $I$ thuộc trục hoành. Hoành độ của điểm $I$ bằng bao nhiêu? \shortans[4]{$4$}
	\loigiai{ Tâm $I$ của đường tròn thuộc trục hoành nên $I(a;0)$.\\
		Đường tròn tâm $I(a;0)$ đi qua hai điểm $A( 1;1)$, $B(5;3)$ nên
		\begin{eqnarray*}
			& & IA=IB\\
			&\Leftrightarrow & \sqrt{(1-a)^2+1^2}=\sqrt{(5-a)^2+3^2}\\
			&\Leftrightarrow & 1-2a+a^2+1=25-10a+a^2+9\\
			&\Leftrightarrow & 8a =32\\
			&\Leftrightarrow & a=4.
		\end{eqnarray*}
		Do đó tâm của đường tròn là $I(4;0)$.
	}
\end{ex}

\Closesolutionfile{ans}

\TL
\begin{ex}
	Khai triển nhị thức $(3x-2)^5$.
\end{ex}

\begin{ex}%[0H9V4-1]
Cho điểm $A(4;2)$ và hai đường thẳng $d\colon x-3y-2=0$ và $\Delta\colon x-3y+18=0$.
\begin{enumerate}
	\item Viết phương trình đường thẳng $d'$ đi qua điểm $A$ và song song với đường thẳng $d$.
	\item Viết phương trình đường tròn $(C)$ đi qua điểm $A$ và tiếp xúc với cả hai đường thẳng $d$ và $\Delta$.
\end{enumerate}
\loigiai{
Xét $d\colon x-3y-2=0$ và $\Delta\colon x-3y+18=0$ có $d\parallel \Delta$. \\
Gọi $D$ là đường thẳng song song và cách đều $d$ và $\Delta$ có $D\colon x-3y+8=0$. \\
Gọi $I$ là tâm của đường tròn $(C)$ ta có $I\in D$ $\Rightarrow I(3b-8;b)$. \\
Ta có $IA=\mathrm{\,d}(I,\Delta) \Leftrightarrow \sqrt{(3b-12)^2+(b-2)^2}=\dfrac{|3b-8-3b+18|}{\sqrt{1^2+(-3)^2}}$
$\Leftrightarrow \hoac{&b=3 \\ &b=\dfrac{23}{5}.}$ \\
\begin{itemize}
\item $b=3 \Rightarrow I(1;3) \Rightarrow R=IA=\sqrt{10}$.\\
Phương trình $(C)\colon (x-1)^2+(y-3)^2=10$.
\item $b=\dfrac{23}{5} \Rightarrow I\left(\dfrac{29}{5};\dfrac{23}{5}\right) \Rightarrow R=IA=\sqrt{10}$.\\
Phương trình $(C)\colon \left(x-\dfrac{29}{5}\right)^2+\left(y-\dfrac{23}{5}\right)^2=10$.
\end{itemize}
}
\end{ex}

\begin{ex}%[Đề chuẩn hóa Toán 10]%[Nguyễn Trung Kiên]%[0D0C2-3]
	Xếp ngẫu nhiên $10$ học sinh gồm $2$ học sinh lớp 10A, $3$ học sinh lớp 10B và $5$ học sinh lớp 10C thành một hàng ngang. Tính xác suất để trong $10$ học sinh trên không có $2$ học sinh nào cùng lớp đứng cạnh nhau.
	% \shortans[]{$630$}
	\loigiai{
	Phép thử \lq\lq Xếp ngẫu nhiên $10$ học sinh gồm $2$ học sinh lớp 10A, $3$ học sinh lớp 10B và $5$ học sinh lớp 10C thành một hàng ngang\rq\rq\, có số phần tử của không gian mẫu là $n(\Omega)=10!$.\\
	Gọi $A$ là biến cố \lq\lq Trong $10$ học sinh trên không có $2$ học sinh cùng lớp đứng cạnh nhau\rq\rq.
	\begin{itemize}
	\item Xếp $5$ học sinh lớp 10C thành một hàng ngang có $5!$ cách. Khi đó với mỗi cách xếp $5$ học sinh lớp 10C sẽ có $6$ khoảng trống gồm $4$ khoảng trống giữa các học sinh lớp 10C và $2$ khoảng trống ở hai đầu.
	\item Xếp $3$ học sinh lớp 10B vào các khoảng trống, sau đó xếp $2$ học sinh lớp 10A vào các khoảng trống còn lại. Dễ thấy không thể xếp đồng thời $2$ học sinh lớp 10B vào $2$ vị trí hai đầu vì khi đó chắc chắn sẽ có ít nhất $2$ học sinh lớp 10C đứng cạnh nhau.\\
	Ta xét các trường hợp sau:
	\begin{itemize}
	\item[TH1:] Không có học sinh 10B ngồi ở vị trí hai đầu hàng của nhóm học sinh lớp 10C đã xếp. Khi đó, xếp $3$ học sinh lớp 10B vào $4$ vị trí trống ở giữa có $\mathrm{A}_4^3$ cách. Ứng với mỗi cách xếp đó, chọn lấy $1$ trong $2$ học sinh lớp 10A xếp vào vị trí trống thứ $4$ (để hai học sinh lớp 10C không được ngồi cạnh nhau), sau đó xếp học sinh lớp 10A còn lại vào $1$ trong $8$ vị trí (không cạnh học sinh lớp 10A đã xếp).\\
	Trường hợp này có $\mathrm{A}_4^3 \cdot 2 \cdot 8 = 384$ khả năng.
	\item[TH2:] Có $1$ học sinh 10B ngồi một trong hai đầu hàng của nhóm học sinh lớp 10C. Khi đó còn $4$ học sinh còn lại ta xếp vào $4$ khoảng trống giữa các học sinh lớp 11C.\\
	Trường hợp này có $\mathrm{C}_3^1 \cdot 2\cdot 4! = 144$ khả năng.
	\end{itemize}
	\end{itemize}
	Vậy số các kết quả thuận lợi của biến cố $A$ là $n(A)=5!\cdot (384+144)= 63\,360$.\\
	Xác suất của biến cố $A$ là $\mathrm{P}(A)=\dfrac{n(A)}{n(\Omega)}=\dfrac{63\,360}{10!} = \dfrac{11}{630}$.
	}
	\end{ex}
\Closesolutionfile{ansbook}

